\documentclass[12pt,a4paper]{article}
\usepackage[UTF8]{ctex}
\usepackage{amsmath, amssymb}
\usepackage{geometry}
\usepackage{booktabs}
\usepackage{enumitem}
\usepackage{tcolorbox}
\usepackage{float}
\usepackage{hyperref}

\geometry{left=2.5cm, right=2.5cm, top=2.5cm, bottom=2.5cm}

\title{\textbf{问题一 模型建立与求解（论文手工作稿）}\\[6pt]
\large VLSI 布图规划设计 --- 自由轮廓面积最小化（面积 + 长宽比）}
\author{2026 年第七届"华数杯"大学生数学建模竞赛 B 题}
\date{2026/08}

\begin{document}
\maketitle
\tableofcontents
\newpage

% ============================================================
\section{问题重述与模型建立}
% ============================================================

\subsection{问题重述}

VLSI 布图规划问题 1 要求在\textbf{芯片轮廓不固定}、\textbf{模块之间不存在
连接关系}的条件下，仅依据各芯片 \texttt{.blocks} 文件中的矩形模块尺寸，
确定每个模块的摆放位置与朝向，使得所有模块外接包围盒的面积最小，并在
面积相同的方案中令包围盒长宽比尽可能接近 1。

\subsection{问题分析}

该问题属于二维矩形装箱与无轮廓布图规划的组合优化问题。搜索空间由模块
排列顺序与朝向组合构成，规模随模块数呈组合爆炸增长：$n$ 个模块的排列
数 $n!$、朝向组合 $2^n$、二叉树形状数为 Catalan 数
$C_n = \frac{1}{n+1}\binom{2n}{n} \sim \frac{4^n}{n^{3/2}\sqrt{\pi}}$，
以 $n=300$ 计仅形状数即约 $10^{176}$。因此必须借助结构化的编码表示将
几何约束内化到搜索空间之中。

本问采用\textbf{基于 B$^*$ 树编码的布图优化模型}，并以\textbf{快速模拟
退火}（Fast-SA）作为求解算法。B$^*$ 树将"任意两模块不重叠且布局紧凑"
的几何要求转化为编码结构的固有性质：给定一棵合法的 B$^*$ 树与各模块
朝向，经过轮廓压缩解码后得到的布局\textbf{必然无重叠}且为\textbf{可容许
布局}（所有模块均无法再向左或向下移动），从而无需在优化过程中显式处理
$O(n^2)$ 对非重叠约束，把"连续坐标 $+$ 几何判定"的难优化问题转化为
有限组合空间上的搜索问题。该框架与题面参考文献 [2] 中的经典方法一致，
适用于 100--300 模块规模，且其编码、解码、扰动与核验环节与目标函数
\textbf{解耦}，便于后续问题仅替换代价函数即可复用。

相较而言，若将非重叠关系写成 0--1 变量的混合整数规划，在模块数达 300
时需引入 $O(n^2)$ 对不重叠约束与 $n$ 组连续坐标变量，约束规模过大，
竞赛时间内难以获得可行解；遗传算法对 B$^*$ 树等装箱结构的交叉维护成本
亦高于模拟退火的局部扰动。故本问采用 B$^*$ 树 + Fast-SA 方案。

\subsection{决策变量与派生量}

设某芯片的模块集合为 $I = \{1, 2, \dots, n\}$，模块 $i$ 的原方向宽、高
与面积分别记为 $w_i$、$h_i$ 与 $a_i = w_i h_i$，总模块面积为
$T = \sum_{i \in I} a_i$。

决策变量包括两大部分：
\begin{enumerate}
    \item \textbf{拓扑变量}：$n$ 个节点的有序二叉树（B$^*$ 树）$\tau$，
    其节点与模块一一对应，隐式编码了全部模块的相对位置关系；
    \item \textbf{朝向变量}：每个模块的旋转标志 $r_i \in \{0, 1\}$，
    $r_i = 0$ 表示保持原方向，$r_i = 1$ 表示旋转 $90^\circ$。旋转后模块
    的有效尺寸为
    \begin{equation}
    (w'_i, h'_i) = \begin{cases}
        (w_i, h_i), & r_i = 0, \\
        (h_i, w_i), & r_i = 1.
    \end{cases}
    \end{equation}
\end{enumerate}

各模块左下角坐标 $(x_i, y_i)$ \textbf{不参与独立搜索}，而是由 $(\tau, r)$
解码唯一确定，属于\textbf{派生量}（见下节）。

\subsection{B$^*$ 树编码与轮廓解码}

B$^*$ 树解码按\textbf{先序遍历}进行：
\begin{enumerate}
    \item 根节点对应的模块置于左下角 $(0, 0)$；
    \item 若节点 $j$ 是节点 $i$ 的\textbf{左孩子}，则模块 $j$ 放置在模块
    $i$ 的右侧且与之相邻，即
    \begin{equation}
    x_j = x_i + w'_i;
    \end{equation}
    \item 若节点 $j$ 是节点 $i$ 的\textbf{右孩子}，则模块 $j$ 与模块 $i$
    左对齐并放置在其上方，即
    \begin{equation}
    x_j = x_i;
    \end{equation}
    \item 纵坐标由\textbf{轮廓数据结构}确定：对候选横区间
    $[x_j,\ x_j + w'_j]$，取其上方已放置模块轮廓高度的最大值作为 $y_j$，
    放置后更新轮廓。
\end{enumerate}

该过程保证每个模块都紧贴左侧或下方的已放置模块，得到所有模块均无法再
向左或向下移动的\textbf{可容许布局}；任意两模块由树结构与轮廓放置唯一
确定而互不重叠，且全部坐标非负。解码结束后，包围盒尺寸为
\begin{equation}
W = \max_{i \in I}\, (x_i + w'_i), \qquad
H = \max_{i \in I}\, (y_i + h'_i).
\end{equation}

\textbf{实现细节}：轮廓线用链表维护，每段至多被覆盖一次，摊还复杂度
$O(n)$；树节点指针采用 10-bit 位压缩（父子/左右子/旋转位共存于一个
32 位整数），支持至 1023 节点。

\subsection{目标函数}

优化目标分为两级：
\begin{itemize}
    \item \textbf{主目标}：包围盒面积 $A = W \times H$ 最小；
    \item \textbf{次级目标}：长宽比 $R = \dfrac{\max(W,H)}{\min(W,H)} \ge 1$
    尽可能接近 1。
\end{itemize}

为保证面积差异优先于长宽比差异、且两项量纲一致，采用加权综合目标：
\begin{equation}
\label{eq:cost}
f(\tau, r) = \lambda\, \frac{A}{\bar A}
           + (1 - \lambda)\, \frac{P(R)}{\bar P},
\end{equation}
其中长宽比惩罚取\textbf{对称形式}
\begin{equation}
P(R) = R + \frac{1}{R} - 2.
\end{equation}

\begin{itemize}
    \item $P(1) = 0$ 且 $P'(1) = 0$：正方形时无惩罚，且附近一阶光滑，
    不过度强迫方形而牺牲面积；
    \item 旋转整个芯片时 $R \to 1/R$，而 $R + 1/R - 2$ 不变：惩罚对称，
    两种等价朝向的布图代价相同。
\end{itemize}

归一化常数 $\bar A$、$\bar P$ 在退火构造期采样 $N+1$ 个随机布局估计：
\begin{equation}
\bar A = \frac{1.1}{N+1} \sum_{k=0}^{N} A_k, \qquad
\bar P = \frac{1.1}{N+1} \sum_{k=0}^{N} P(R_k),
\end{equation}
系数 $1.1$ 为松弛因子，防止采样均值过低导致权重失衡。两项除以统计均值
后均为 $O(1)$ 量级，$\lambda$ 因此具有"面积与长宽比的相对重要性"的直观
含义。

\textbf{权重 $\lambda$ 的确定}：扫描 $\lambda \in [0.45, 0.95]$ 得到权衡
曲线（见 \ref{sec:lambda}）：$\lambda = 0.8$ 时退火陷入长条局部最优
（n300 长宽比高达 4.4）；$\lambda = 0.5$ 时 n100 长宽比可达 1.002、
面积利用率 0.942。$\lambda = 0.5$ 为面积与长宽比的均衡点，故取
$\lambda = 0.5$。

\textbf{模型合理性}：
\begin{itemize}
    \item 硬约束（不重叠、尺寸保持）由解码器自动满足，目标函数无需罚项，
    不存在罚权重调参问题；
    \item 面积下界 $WH \ge T$ 恒成立，故面积利用率
    $\eta = T/(WH) \le 1$ 是无偏的解质量度量（$\eta = 1$ 仅对极特殊尺寸
    集合可达）。
\end{itemize}

% ============================================================
\section{求解算法设计}
% ============================================================

\subsection{总体框架}

\begin{tcolorbox}[colback=blue!5, colframe=blue!60, title=算法流程]
\begin{enumerate}
    \item 用 \textbf{B$^*$ 树}表示布局拓扑（每个模块是树的一个节点）；
    \item 用 \textbf{轮廓压缩解码}将树确定性解码为无重叠可容许布图；
    \item 用 \textbf{快速模拟退火}（Fast-SA）在 B$^*$ 树空间搜索：
    每步执行一次扰动算子、重新解码、按式 \eqref{eq:cost} 计算代价；
    \item 每芯片独立运行多轮，按字典序 $(A, R)$ 取最优（面积优先，
    长宽比次之）。
\end{enumerate}
\end{tcolorbox}

\subsection{扰动算子}

退火每步以一定概率执行三种算子之一，形成完备邻域：

\begin{table}[H]
\centering
\small
\caption{三种扰动算子及其概率}
\begin{tabular}{clc}
\toprule
算子 & 操作语义 & 概率 \\
\midrule
旋转 (rotate) & 随机选一模块翻转 $90^\circ$（仅翻转尺寸标志） & 0.30 \\
删除-插入 (del\_and\_ins) & 随机删一节点再插入到随机位置（改变父子结构） & 0.35 \\
节点交换 (swap) & 随机交换两节点（父子相邻走近邻交换，否则远邻交换） & 0.35 \\
\bottomrule
\end{tabular}
\end{table}

\begin{itemize}
    \item \textbf{局部性互补}：删除-插入只改变局部父子结构，解码后多只
    影响局部区域，利于精细搜索；节点交换（远邻）跳跃更大，利于跳出
    局部最优；
    \item \textbf{拓扑保持}：删除-插入与交换均保持树拓扑合法（节点集合
    守恒、无环），由 \texttt{Validate\-Tree()} 白盒校验。
\end{itemize}

\subsection{快速模拟退火}

\subsubsection{构造期：归一化常数与初始温度}

构造期执行 $N+1$ 次随机扰动 $+$ 解码，统计式 \eqref{eq:cost} 的分母
$\bar A$、$\bar P$，并统计相邻代价差的平均幅度 $\bar{|\Delta|}$，用于
估计初始温度，使初始接受率约为 $P = 0.9$：
\begin{equation}
T_0 = -\,\frac{\bar{|\Delta|}}{\ln P} \approx 9.49\, \bar{|\Delta|}.
\end{equation}

\subsubsection{自适应冷却}

\textbf{不用固定几何冷却}（$T \leftarrow \alpha T$）的原因：其收敛速度对
问题规模与代价尺度敏感，参数 $\alpha$ 难以通用。本项目采用\textbf{平均
增量驱动}的自适应降温：每一温度层执行 $rnd = 2n + 20$ 次扰动，统计层内
平均代价差 $\bar{|\Delta|}_k$，按式更新：
\begin{equation}
T_k = \begin{cases}
T_0 \cdot \dfrac{\bar{|\Delta|}_k}{c \cdot k}, & k \le k_0, \\[1.2em]
T_0 \cdot \dfrac{\bar{|\Delta|}_k}{k}, & k > k_0,
\end{cases}
\qquad
k_0 = \max\!\Big(2,\ \Big\lceil \tfrac{n}{11} \Big\rceil\Big), \quad
c = \max(100 - n,\ 10).
\end{equation}

物理含义：层内平均代价差大（搜索激烈）时降温快，趋于平稳时降温慢，
实现"粗段快冷、细段慢冷"的自动平衡。

\subsubsection{接受准则（Metropolis）}

对候选布图代价差 $\Delta = f_{\text{new}} - f_{\text{cur}}$：
\begin{equation}
\text{接受} \iff \Delta \le 0 \quad \text{或} \quad
\xi < \exp\!\Big(-\frac{\Delta}{T}\Big), \qquad \xi \sim U(0,1).
\end{equation}
高温期接受劣解概率高（全局探索），低温期收敛于局部精细搜索。被接受时
更新当前解与历史最优；被拒绝时回滚拓扑。

\subsubsection{停滞重置（reheating）}

连续 $2n$ 个温度层无改进时判定为停滞，重置 $T \leftarrow T_0$ 并递增
$rnd$（扩大邻域规模），重新加热搜索；累计重置超过 $n/7$ 次（精修阶段）
则终止。该机制可多次跳出深层局部最优。

\subsubsection{终止条件}

\begin{table}[H]
\centering
\small
\begin{tabular}{ll}
\toprule
条件 & 含义 \\
\midrule
$T \le 0.001$ & 温度降到下界（冻结） \\
拒绝率 $\ge 99\%$ & 连续 $rnd$ 次几乎全部拒绝（收敛） \\
$k > 9n$（精修） & 迭代层数上限 \\
停滞重置超限 & 重置次数 $> n/7 + 1$ \\
\bottomrule
\end{tabular}
\end{table}

\subsubsection{整体伪代码}

\begin{tcolorbox}[colback=blue!5, colframe=blue!60, title=算法 1：B$^*$ 树 + 快速模拟退火]
{\small
\begin{verbatim}
输入: 模块尺寸, lambda = 0.5, 采样次数 N+1
输出: 最优布图 (坐标与旋转)

1  采样 N+1 个随机布局 -> 归一化常数 A_bar, P_bar, 初始温度 T0
2  T <- T0/50; 初始化 B*-Tree
3  while 未满足终止条件 do
4      for i <- 1 to rnd (= 2n+20) do
5          tau' <- 扰动(tau)                # 旋转 / 删除-插入 / 交换
6          轮廓解码 tau' -> (W, H); 计算代价 f(tau')
7          d <- f(tau') - f(tau)
8          if d <= 0 或 xi < exp(-d/T) then 接受 tau'
9      T <- T0 * avg(|d|) / (c * k)         # 自适应降温
10     if 停滞 then T <- T0                 # reheating
11 return 历史最优布图
\end{verbatim}
}
\end{tcolorbox}

说明：精修阶段初始温度取 $T_0 / 50$，从较低温度开始聚焦局部最优；
每芯片执行两轮独立精修，取代价更小者输出。

\subsection{多轮独立运行取优}

模拟退火是随机启发式，单次运行结果有波动（实测 n300 长宽比在
1.5 $\sim$ 4.4 间波动）。因此每芯片独立运行 $R$ 轮（交付结果取 $R = 5$，
默认 $R = 3$，可用 \texttt{--repeats} 调整），按字典序 $(A, R)$ 取最优，
将随机波动压到可接受范围。

\subsection{复杂度分析}

\begin{table}[H]
\centering
\small
\begin{tabular}{lc}
\toprule
环节 & 复杂度 \\
\midrule
单次解码 & $O(n)$（摊还） \\
单温度层 & $O(rnd \cdot n) = O(n^2)$，$rnd = 2n + 20$ \\
单轮退火 & 温度层数 $\times$ $O(n^2)$，层数由降温/重置动态决定 \\
单芯片 $R$ 轮 & $R \times$ 单轮（三芯片可并行） \\
\bottomrule
\end{tabular}
\end{table}

实测单轮耗时：n100 $\approx$ 0.17s、n200 $\approx$ 3.2s、n300
$\approx$ 8.8s，随 $n$ 近似线性扩展（$n$ 翻倍约 3 倍），工程可行。

% ============================================================
\section{结果与验证}
% ============================================================

\subsection{数据规模}

\begin{table}[H]
\centering
\caption{三组芯片数据规模}
\begin{tabular}{lccc}
\toprule
芯片 & 模块数 $n$ & 端子数 & 线网数 \\
\midrule
n100 & 100 & 334 & 885 \\
n200 & 200 & 564 & 1585 \\
n300 & 300 & 569 & 1893 \\
\bottomrule
\end{tabular}
\end{table}

\subsection{求解结果}

% TODO: λ 扫描完成后，以 q1/output/q1_metrics.csv 与 data/raw 重算
% 下表（面积利用率 = T/A、死区比例 = 1 - T/A、下界相对差距 = (A-T)/T、
% 宽高比 R = max(W,H)/min(W,H)）。

\begin{table}[H]
\centering
\caption{Q1 三芯片求解结果（$\lambda = 0.5$，5 轮取优）}
\small
\begin{tabular}{lcccccccc}
\toprule
芯片 & 轮廓 $W \times H$ & 面积 $A$ & 模块总面积 $T$ & 面积利用率 $\eta$ & 死区比例 & 下界相对差距 & 宽高比 $R$ & 耗时(s) \\
\midrule
n100 & $421\times455$ & 191555 & 179501 & 0.9371 & 6.29\% & 6.72\% & 1.081 & 0.85 \\
n200 & $357\times517$ & 184569 & 175696 & 0.9519 & 4.81\% & 5.05\% & 1.448 & 15.9 \\
n300 & $486\times591$ & 287226 & 273170 & 0.9511 & 4.89\% & 5.15\% & 1.216 & 43.9 \\
\bottomrule
\end{tabular}
\end{table}

注：面积利用率 $\eta = T/A$（$T$ 为模块总面积，面积下界）；死区比例
$= 1 - \eta$；下界相对差距 $= (A - T)/T$。三组芯片均通过独立合法性复核
（无重叠、尺寸保持）。

\subsection{结果讨论}

\begin{itemize}
    \item \textbf{面积利用率与死区}：三组芯片利用率均 $> 93\%$（死区比例
    $< 7\%$）。理论上界 $\eta = 1$ 仅对极特殊尺寸集合可达，故 $\eta > 0.93$
    已接近工程极限；剩余死区面积 n100 / n200 / n300 分别为 12054 / 8873 /
    14056，占总面积的 6.29\% / 4.81\% / 4.89\%；
    \item \textbf{长宽比}：全部 $\le 1.5$，满足竞赛验证标准；n100 最接近
    正方形（1.081），n200 最长条（1.448）；
    \item \textbf{规模扩展性}：耗时随 $n$ 近似线性增长，说明自适应降温
    与 reheating 参数对规模不敏感，算法可迁移到更大规模实例；
    \item \textbf{字典序取优}：面积优先、长宽比次之，二者未发现冲突
    （面积最小的解长宽比也可控）。
\end{itemize}

\subsection{权重 $\lambda$ 敏感性}
\label{sec:lambda}

\begin{table}[H]
\centering
\small
\begin{tabular}{ll}
\toprule
$\lambda$ & 观察结果 \\
\midrule
0.5 & n100 长宽比可达 1.002、利用率 0.942（均衡点，交付采用） \\
0.8 & n300 卡长条局部最优，长宽比高达 4.4（面积权重过大） \\
$\to 1$ & 完全忽略长宽比，退化为纯面积最小化，长宽比不可控 \\
\bottomrule
\end{tabular}
\end{table}

\subsection{独立合法性复核}

求解器输出坐标后，用\textbf{独立于 C++ 的 Python 实现}复核（避免实现
耦合导致的系统性错误）：
\begin{enumerate}
    \item 模块个数与输入一致；
    \item 两两无重叠（严格相交语义）；
    \item 每模块有效尺寸恰为原始尺寸的置换（旋转不改变尺寸集）；
    \item $WH \ge T$（面积下界）。
\end{enumerate}
三组芯片全部通过。

% ============================================================
\section{正确性保障与可复现性}
% ============================================================

\subsection{工程保障}

\begin{itemize}
    \item 核心代码正确性由 \textbf{1027 断言白盒测试} + ASan/UBSan 双轨
    保证：含独立参考解码器（与实现完全不同的算法）、暴力枚举全部二叉树
    形状（Catalan 132/429/1430）、独立代价 oracle、算子语义校验；
    \item 求解与可视化解耦：求解只写文件，画图后处理读取，不互相拖累；
    \item 原始数据 \texttt{data/raw/} 只读，求解器不修改任何输入文件；
    \item 随机种子未固定（每次运行结果略有差异），论文数值以
    \texttt{q1\_metrics.csv} 为准，复现流程（命令）确定。
\end{itemize}

\subsection{运行命令与产物}

{\small
\begin{verbatim}
cd cpp_solver && make && make test      # 编译 + 白盒测试（1027 断言）
conda activate math
python q1/q1_solver.py                  # 求解三芯片 + 汇总 CSV + 出图
\end{verbatim}
}

\begin{itemize}
    \item \texttt{q1/output/q1\_metrics.csv}：三芯片指标汇总（论文数值以此为准）；
    \item \texttt{q1/output/q1\_<chip>.rpt}：最终布局坐标（格式：
    cost / hpwl / area / W H / time / 每行 \texttt{name x1 y1 x2 y2}）；
    \item \texttt{q1/output/q1\_<chip>.log}：收敛轨迹 + 9 等分过程快照；
    \item \texttt{q1/visualization/figs/<时间戳>/<chip>/}：每芯片
    9 张过程快照 + 1 张收敛曲线 + 1 张最终布图（构成收敛性的可视化证据）。
\end{itemize}

% ============================================================
\begin{thebibliography}{9}
\bibitem{bstar} Chang Y C, Chang Y W, Wu G M, et al. B*-Trees: a new
  representation for non-slicing floorplans[C]// Proceedings of the 37th
  Annual Design Automation Conference (DAC). 2000: 458--463.
\bibitem{fastsa} Chen T C, Chang Y W. Modern floorplanning based on B*-tree
  and fast simulated annealing[J]. IEEE Transactions on Computer-Aided
  Design of Integrated Circuits and Systems, 2006, 25(4): 637--650.
\end{thebibliography}

\end{document}
